% =====================================================================
%  Chapter 1 — Why? and Why Not?
%  Map: research/chapter-map.md (Chapter 1)
%  Rule: the author is never quoted in body text; short quotes only in
%  footnotes. Full statements live in interviews/.
% =====================================================================
\chapter{Why? and Why Not?}
\label{ch:why}

% TODO-INTERVIEW: Choose an epigraph (not from the author's own words).

Every programming language is an answer. Most of the time nobody remembers
the question. Garbage collectors, loops, exceptions, type aliases, optimizers
that quietly repair what the programmer wrote: these arrive in a new language
as inheritance, carried forward because the last language had them and the one
before that did too. Silica began when Lee Barney stopped accepting the
inheritance and started asking, of each piece, a child's question: \emph{why?}

This book is about what happened when that question was asked persistently,
for years, of the tools working programmers trust most. It is also about the
question that comes second, and is harder to ask: \emph{why not?} The first
question takes things apart. The second decides whether to build something in
their place.

\section{The question that would not go away}

Barney cannot point to the first time the question came up. What stands out in
memory is the arrival of Java, and how much safer it felt than the C and C++
that Barney had been writing in industry. The sense that some languages protect
the people who use them, and others leave them exposed, was there long before
it had a name.

In December 2023 a group of national security agencies published a case for
memory-safe roadmaps, and in early 2024 the White House technical report
followed. Both named a short list of languages---C\#, Go, Java, Python, Rust and
Swift---as memory safe. Reading them, Barney, then teaching at Brigham Young
University--Idaho, wondered how that list had been chosen. Why those languages,
and not others?

The paper that followed was not Barney's idea. At CodeBEAM in San Francisco, a
professional conference for the Erlang community rather than an academic one,
a member of the Erlang Ecosystem Foundation asked Barney to write an article
arguing, if the data supported it, that the BEAM languages Erlang and Elixir
belonged on the list.\footnote{Interview, 2026-11-13.}

The data did not entirely cooperate. The paper, written with Adolfo Neto and
presented at the ACM Erlang Workshop in Milan in September
2024,\footnote{L.~Barney and A.~Neto, ``Unsafe Impedance: Safe Languages and
Safe by Design Software,'' Erlang '24. Source:
\src{BEAM Security Article/Article/unsafe_impedance_0715/}.} argued that being
memory safe is a necessary property and not a sufficient
one.\footnote{The paper's phrase is ``necessary but insufficient.''} It proposed
a new way to measure languages: by how hard they make it to write, load, and
above all \emph{hide} unsafe code. The paper called that difficulty
\emph{unsafe impedance}, and graded languages from none to infinite. Erlang and
Elixir scored well, but not perfectly, because native code loaded into the BEAM
can hide too. And the paper's last observation was the one that mattered. No
language that was secure by default had infinite unsafe impedance. None
existed.\footnote{``there are no secure by default languages that have
infinite unsafe impedance.''}

An advocate might have been troubled by where the evidence led. Barney was not.
The point, as Barney saw it, was to find out what was true rather than to win
an argument, and the truth was that no language, the BEAM's included, was beyond
criticism. The same standard would later be turned on Silica
itself.\footnote{``Advocacy is not what I'm about. I care about honesty and
truth.'' Interview, 2026-11-13.}

That last sentence of the paper is a strange one to write if you have no
intention of doing anything about it. Chapter~\ref{ch:danger} follows the idea
into Silica's foreign function interface. Here it matters for a different
reason: it is the first place in the record where the \emph{why?} points at a
gap rather than at a fault.

\section{A habit of asking}

Read Silica's design documents in order and the same question turns up again
and again, in different clothes.

Why does a language need a garbage collector? Silica has none, and its memory
model (Chapter~\ref{ch:memory}) is built to make one unnecessary.
Why does a language need loops, when recursion can say the same thing in one
form instead of three? Why let a programmer give a type a new name, when the
name hides the shape a reader needs to see? Why let the optimizer clean up dead
code and duplicate work, rather than asking the programmer to write what they
meant? Why do languages still model the machine as it was in the 1970s?
Why should an idle program end itself, or a runtime guess at defaults the
program never chose?

Each of these questions has a conventional answer, and the conventional
answers are not foolish. Garbage collection frees programmers from a whole
class of error. Optimizers make naive code fast. Defaults make simple programs
short. What the questions expose is that each convenience comes with a cost
that someone, somewhere, stops being able to see. Silica's principles,
the subject of the second part of this book, are largely a list of costs
made visible again.

% TODO-INTERVIEW: Which "why?" was the most uncomfortable to ask — the one
%   where the conventional answer turned out to be strongest?

\section{When why becomes why not}

Asking why is cheap. It is possible to spend a career asking it, writing
papers, and changing nothing. Barney's record from 2023 to 2025 shows four
attempts to answer the question \emph{inside} existing systems before the
decision to step outside them.

The first was SwErl, a library that brought Erlang's processes and its
proven design patterns into Swift (Chapter~\ref{ch:history}). The second
was the 2024 paper, which proposed a human review process, the Unsafe
Acceptance Process, to raise a team's unsafe impedance when the language would
not. The third was BEAM\_SEMP, a secure overlay for Erlang networking, which
spent the autumn of 2025 pulling dangerous defaults out of a running BEAM one
at a time: switching off standard distribution, banning remote atom creation
and code loading, refusing to tell a rejected peer why.

The fourth, running alongside the others through 2025, was the most ambitious:
a port of the BEAM itself from C to Rust. Barney had grown tired of telling
audiences that Erlang was a memory-safe language, only to be asked what its
virtual machine was written in, and having to answer, truthfully, C. The
inevitable reply was that it could not really be safe, then. A virtual machine
rebuilt in a memory-safe language would answer the objection.

Each attempt worked, after a fashion. Each also ran into the same ceiling.
A library cannot fix the language it is written in. A review process depends on
people never having a bad day. A security layer bolted onto a permissive
runtime is a list of patches, and the next default nobody thought of is
waiting underneath. And a port, however faithful, can only carry an old design
into a new language.

% TODO-VERIFY: BEAM_SEMP's last commit date: 2025-10-09 or 2025-10-31.
% TODO-INTERVIEW: The Rust BEAM port — where is it now, and how far did it get?

\section{The wrong problem}

By 24~December 2025 the Rust port was well along; the re-engineering of the
virtual machine was most of the way to completion. Barney was working on it
that day, paused, and began to think. What came out of the thinking was that the
port was solving the wrong problem.

Every language is designed to solve a particular set of problems. Rust was not
designed to express the ideas and principles at the heart of the BEAM, and
making it do so meant forcing square pegs into round holes. No amount of effort
would produce what an engineer would call a clean solution, and the same was
true of every other language Barney knew. Barney spent a few minutes going
through them anyway, looking for one that could readily express the ideas of
SwErl and SEMP and still be memory safe. It took only those few minutes to see
that there was none.

The first designs for a new language, on paper and in Barney's head, took a few
hours. Two days later, on 26~December 2025, Silica's first commit was made. At
that moment the question had stopped being \emph{why is it like this?} and
become \emph{why not build something that is not?}

% TODO-INTERVIEW: The 24 December scene — where were you, and what was on the
%   screen when you paused?

\section{A recovery, and a new kind of reader}

The \emph{why not?} deserves a fuller answer, because the person asking it in
December 2025 had spent the previous six months relearning how to think.

In June 2025 Barney had what is commonly called a mini-stroke. For a time
Barney could not read. Thinking was difficult; thinking about large and complex
things was impossible; forming sentences in order to speak was severely
restricted. The work on the Rust port became part of the recovery. Rebuilding a
virtual machine, piece by piece, helped bring back the ability to hold large
structures in mind, and it was that regained ability that eventually made it
possible to see the port as the wrong problem. Barney tells this part of the
story openly, so that others who face something similar know that they are not
alone.\footnote{Interview, 2026-11-13.}

During the recovery, language models did much of the coding that Barney could
not. Watching them work, Barney noticed how badly they struggled with existing
languages, all of which had been designed to be understood by humans alone.
Getting the code Barney wanted out of them meant fighting them. It began to seem
easier to start from nothing: to build a language and a compiler around what
language models do well, instead of forcing them to do what they do badly.
That observation, made out of necessity, became one of Silica's founding
principles (Chapter~\ref{ch:oneway}).

None of this would surprise anyone who knows Barney. A lifelong lover of
languages of every kind, Barney once considered majoring in linguistics as an
undergraduate, and has long held that language is language, whether spoken by
people or compiled for machines: each one is shaped by the problems it was made
to solve. The same recovery has since carried Barney into learning Tagalog.

\section{Turning the question inward}

There is one more thing the record shows, and it may be the most important.
The questions did not stop once Silica existed. They were turned on Silica
itself.

Generics were removed six days after the first commit. Structs came and went.
The standard data structures were built, judged too complicated for their
purpose, and deleted wholesale in June 2026. Networking was made part of the
language in September and taken out again within days. And in October 2026, in
the middle of designing a new memory model, the project discovered that its
runtime had been quietly using a per-actor heap that nobody had ever
decided to build---while the public documentation promised that memory lived on
actor stacks. The question \emph{why is this here?} started the largest
redesign in the project's history (Chapter~\ref{ch:memory}).

A language designed by asking why has to be willing to ask it of its own
decisions. The reversals in Silica's history are not signs of indecision. They
are the method working.

% TODO-INTERVIEW: Which reversal was the hardest to make? Which "why not?"
%   decision surprised you most once you'd made it?

\section{What the questions became}

Near the end of a 2026 talk introducing Silica to an Erlang audience, the
closing slide posed three questions instead of conclusions: what a compiler
should refuse, where a runtime should isolate failure, and how modern hardware
and tools should shape the design of a source language.\footnote{Slide 11 of
\src{silica/Silica_intro_FP_Erlang_Elixir_Gleam.pptx}, which asks the reader to
``Treat systems programming as a design problem again.''} The rest of this book
is an attempt to show how Silica answers them.

Chapter~\ref{ch:history} tells the story of how the language got from a Rust
bootstrap compiler to a self-hosting compiler that runs from a Mac to a
512-kilobyte microcontroller. Chapters~\ref{ch:hidden} through~\ref{ch:transport}
take the ten principles one at a time. Each began as a \emph{why}. Each became
real only after someone asked \emph{why not?}

\begin{sidebox}[The questions]
\begin{itemize}\itemsep0pt
  \item Why is a language called safe if it lets unsafe code hide?
  \item Why a garbage collector? Why a heap at all?
  \item Why loops, generics, type names, conversions?
  \item Why let the optimizer fix what the programmer wrote?
  \item Why model a 1970s machine?
  \item Why defaults the program never chose?
  \item Why compile libraries separately?
  \item Why patch the runtime instead of owning it?
\end{itemize}
\end{sidebox}
